\documentclass[11pt,a4paper]{article}

\usepackage[utf8]{inputenc}
\usepackage[T1]{fontenc}
\usepackage[brazil]{babel}
\usepackage{lmodern}
\usepackage[a4paper, margin=2.5cm]{geometry}
\usepackage{amsmath}
\usepackage{booktabs}
\usepackage{float}
\usepackage{caption}
\usepackage{fancyhdr}
\usepackage{pgfplots}
\pgfplotsset{compat=1.18}
\usepackage[siunitx]{circuitikz}
\usepackage[hidelinks]{hyperref}

\captionsetup{font=small, labelsep=endash}
\setlength{\parindent}{1.25cm}
\setlength{\parskip}{4pt}

\pagestyle{fancy}
\fancyhf{}
\fancyhead[C]{\small Projeto Interdisciplinar -- Engenharia de Computação, 2º sem. 2026 -- IFSP Guarulhos}
\fancyfoot[C]{\thepage}
\renewcommand{\headrulewidth}{0.4pt}

\begin{document}

\begin{center}
	{\Large\bfseries LUXIF: Luxímetro de baixo custo com LDR e faixa automática}

	\vspace{12pt}

	{\small
	EDUARDO L. MARTINEZ\textsuperscript{1}, EMILY O. ARAUJO\textsuperscript{2},
	HENRIQUE F. DEFAVERE\textsuperscript{3}, RAFAEL SHIGUEMI U. RAMOS\textsuperscript{4},
	VINÍCIUS M. PEREIRA\textsuperscript{5}}
\end{center}

\vspace{4pt}
{\footnotesize
\noindent\textsuperscript{1}Graduando em Engenharia de Computação, IFSP, Câmpus Guarulhos, eduardo.martinez@aluno.ifsp.edu.br;\\
\textsuperscript{2}Graduanda em Engenharia de Computação, IFSP, Câmpus Guarulhos, araujo.emily@aluno.ifsp.edu.br;\\
\textsuperscript{3}Graduando em Engenharia de Computação, IFSP, Câmpus Guarulhos, h.defavere@aluno.ifsp.edu.br;\\
\textsuperscript{4}Graduando em Engenharia de Computação, IFSP, Câmpus Guarulhos, rafael.shiguemi@aluno.ifsp.edu.br;\\
% TODO: acrescentar o e-mail institucional do Vinícius (prontuário GU3080315) ao final da linha abaixo
\textsuperscript{5}Graduando em Engenharia de Computação, IFSP, Câmpus Guarulhos}

\vspace{8pt}
\noindent\textbf{Área de conhecimento (Tabela CNPq):} 1.05.00.00-6 Física

\vspace{8pt}
\noindent\textbf{RESUMO:} A medição da intensidade luminosa é uma etapa importante em experimentos de óptica e em aplicações de monitoramento ambiental. Este trabalho apresenta o hardware do LUXIF, um luxímetro de baixo custo baseado em um sensor fotoresistivo (LDR) de 10~mm, em quatro resistores de referência e em um microcontrolador Arduino Uno, com uma versão em ESP32 que transmite as leituras por Wi-Fi para uma página web. O LDR forma um divisor de tensão com um dos resistores, escolhido automaticamente pelo microcontrolador conforme a leitura do conversor analógico-digital, o que amplia a faixa útil de medição e melhora o aproveitamento da sua resolução. A faixa em uso é indicada por LEDs ligados em saídas próprias, fora do caminho da medição, e o valor calculado em lux é exibido em um display OLED. O circuito é alimentado por uma bateria de 9~V ligada ao conector de alimentação da placa. A análise teórica indica que as quatro faixas se sobrepõem e cobrem, juntas, resistências de cerca de 100~$\Omega$ a 4~M$\Omega$, tornando a arquitetura adequada para aplicações didáticas e experimentais.

\vspace{4pt}
\noindent\textbf{PALAVRAS-CHAVE:} fotoresistor; instrumentação eletrônica; divisor de tensão; iluminância; Arduino; ESP32; sensores.

\section*{INTRODUÇÃO}

A medição da intensidade luminosa é fundamental para a realização de experimentos em óptica, permitindo a análise quantitativa de fenômenos como absorção, transmissão e polarização da luz. A utilização de instrumentos de baixo custo para esse fim contribui para a integração entre conceitos teóricos e aplicações práticas em disciplinas de Física e Eletrônica (HALLIDAY; RESNICK; WALKER, 2016).

Os fotoresistores (LDRs) são amplamente utilizados em sistemas de medição luminosa devido à sua simplicidade e ao fato de apresentarem variação de resistência elétrica em função da intensidade de luz incidente (SCHERZ; MONK, 2013). Entretanto, essa resistência varia em várias ordens de grandeza, de centenas de ohms sob luz forte a milhões de ohms no escuro, o que reduz a sensibilidade das medições quando se utiliza apenas um resistor fixo no circuito de leitura.

Diante desse problema, este trabalho propõe um luxímetro em que o microcontrolador escolhe, entre quatro resistores de referência, aquele mais adequado à luz do momento, de forma semelhante à troca automática de escala de um multímetro. O objetivo é projetar um instrumento de baixo custo, simples de montar e capaz de medir diferentes níveis de iluminância, visando sua aplicação em experimentos de óptica e em medições de iluminação natural.

\section*{MATERIAL E MÉTODOS}

O projeto foi desenvolvido a partir de pesquisa bibliográfica em livros, manuais e folhas de dados dos componentes, seguida da modelagem teórica do circuito. Os componentes utilizados estão listados na Tabela~\ref{tab:componentes}.

\begin{table}[H]
	\centering
	\caption{Componentes do LUXIF.}
	\label{tab:componentes}
	\small
	\begin{tabular}{lcl}
		\toprule
		Componente & Qtd. & Função \\
		\midrule
		Arduino Uno & 1 & Trocar as faixas, ler a tensão e calcular o lux \\
		LDR 10 mm & 1 & Sensor de luz \\
		Resistores de 1, 10, 100 e 470 k$\Omega$ & 1 de cada & Resistores de referência das faixas \\
		LED azul & 4 & Indicar a faixa em uso \\
		Resistor de 1 k$\Omega$ & 4 & Limitar a corrente dos LEDs azuis \\
		Display OLED 0,96" (SSD1306, I\textsuperscript{2}C) & 1 & Exibir a leitura \\
		Bateria de 9 V e suporte & 1 & Alimentação portátil \\
		Chave gangorra & 1 & Ligar e desligar o circuito \\
		LED vermelho e resistor de 1 k$\Omega$ & 1 de cada & Indicar que o circuito está ligado \\
		Placa perfurada e fios & -- & Montagem \\
		\bottomrule
	\end{tabular}
\end{table}

\subsection*{Circuito de medição}

O LDR é ligado entre o ponto de leitura e o terra (GND). Cada um dos quatro resistores de referência ($R_n$) é ligado entre esse mesmo ponto e uma saída digital própria do Arduino (D2 a D5). Apenas uma saída é colocada em nível alto (5~V) por vez; as outras são configuradas como entrada, ficando em alta impedância, o que na prática desconecta seus resistores do circuito. O resistor ativo e o LDR formam um divisor de tensão, cuja saída é lida pela entrada analógica A0 (SCHERZ; MONK, 2013):

\begin{equation}
	V = 5 \times \frac{R_{LDR}}{R_n + R_{LDR}}
	\label{eq:divisor}
\end{equation}

\noindent em que $V$ é a tensão em A0, em volts; $R_{LDR}$ é a resistência do LDR e $R_n$ a do resistor ativo, ambas em ohms. Como as duas resistências são positivas, $V$ fica sempre entre 0 e 5~V, que é o limite da entrada analógica. Além disso, mesmo com o menor resistor (1~k$\Omega$), a corrente na saída do Arduino não passa de 5~mA, bem abaixo do limite recomendado de 20~mA por pino.

O conversor analógico-digital do Arduino Uno tem 10 bits e usa os mesmos 5~V como referência, devolvendo uma leitura $N$ entre 0 e 1023. Combinando essa relação com a Equação~\ref{eq:divisor}, a resistência do LDR é obtida diretamente da leitura, sem depender do valor exato da alimentação:

\begin{equation}
	R_{LDR} = R_n \times \frac{N}{1023 - N}
	\label{eq:resistencia}
\end{equation}

\subsection*{Faixa automática}

A cada medição, o microcontrolador lê a entrada A0 dez vezes e calcula a média, para reduzir o ruído. Em seguida, compara o resultado com dois limites:

\begin{itemize}
	\item se $N > 920$ (cerca de 4,5~V), o ambiente está escuro demais para a faixa atual e o próximo resistor, maior, é selecionado;
	\item se $N < 100$ (cerca de 0,5~V), o ambiente está claro demais e o resistor anterior, menor, é selecionado;
	\item caso contrário, a leitura está na região útil e o valor em lux é calculado.
\end{itemize}

Após cada troca, o programa aguarda 10~ms para o circuito estabilizar antes da próxima leitura. A Tabela~\ref{tab:faixas} mostra a resistência do LDR coberta por cada faixa.

\subsection*{Indicação da faixa}

Cada LED azul é ligado, em série com um resistor de 1~k$\Omega$, entre uma saída exclusiva do Arduino (D6 a D9) e o terra, e é aceso junto com a faixa correspondente. Dessa forma, os LEDs ficam fora do divisor de tensão: a queda de tensão do LED azul (de 2,5 a 3~V) não reduz a tensão disponível no divisor nem varia com a corrente, e a medição não é afetada. Com o resistor de 1~k$\Omega$, a corrente em cada LED fica em torno de 2~mA, suficiente para a indicação visual e bem abaixo do limite de 20~mA por pino.

\subsection*{Alimentação}

O circuito é alimentado por uma bateria de 9~V ligada, através de uma chave gangorra, ao conector de alimentação do Arduino Uno, cujo regulador interno aceita de 7 a 12~V e fornece os 5~V usados pelo restante do circuito. O LED vermelho indicador é ligado, com um resistor de 1~k$\Omega$, entre a saída da chave e o terra, em paralelo com a carga, e consome cerca de 7~mA.

O display OLED é ligado ao barramento I\textsuperscript{2}C do Arduino (A4 e A5) e alimentado pelos 5~V e pelo GND da placa. O esquema completo é apresentado na Figura~\ref{fig:esquema}.

\begin{figure}[htbp]
	\centering
	\resizebox{0.95\textwidth}{!}{\begin{circuitikz}[american, scale=0.95, transform shape]
  % Arduino
  \draw[thick] (0,0) rectangle (3,9);
  \node[align=center, font=\bfseries] at (1.5,4.6) {Arduino\\Uno};

  % Pinos do lado direito: faixas e leitura
  \foreach \y/\p in {8.5/D2, 7.5/D3, 6.5/D4, 5.5/D5} {
    \node[anchor=east, font=\small] at (3,\y) {\p};
  }
  \node[anchor=east, font=\small] at (3,3.5) {A0};

  \draw (3,8.5) -- (4,8.5) to[R, l=$1\,\mathrm{k\Omega}$] (6.5,8.5) -- (7.5,8.5);
  \draw (3,7.5) -- (4,7.5) to[R, l=$10\,\mathrm{k\Omega}$] (6.5,7.5) -- (7.5,7.5);
  \draw (3,6.5) -- (4,6.5) to[R, l=$100\,\mathrm{k\Omega}$] (6.5,6.5) -- (7.5,6.5);
  \draw (3,5.5) -- (4,5.5) to[R, l=$470\,\mathrm{k\Omega}$] (6.5,5.5) -- (7.5,5.5);
  \draw (7.5,8.5) -- (7.5,3.5);
  \foreach \y in {7.5, 6.5, 5.5} { \fill (7.5,\y) circle (2pt); }
  \draw (3,3.5) -- (7.5,3.5);
  \fill (7.5,3.5) circle (2pt);
  \draw (7.5,3.5) to[phR, l=LDR] (7.5,1) node[ground]{};

  % LEDs azuis do lado esquerdo
  \foreach \y/\p in {8.5/D6, 7.5/D7, 6.5/D8, 5.5/D9} {
    \node[anchor=west, font=\small] at (0,\y) {\p};
    \draw (0,\y) -- (-0.8,\y) to[led] (-2.4,\y) to[R, l_=$1\,\mathrm{k\Omega}$] (-4.6,\y) -- (-5.2,\y);
  }
  \node[font=\small] at (-1.6,9.3) {LEDs azuis};
  \draw (-5.2,8.5) -- (-5.2,4.8) node[ground]{};
  \foreach \y in {7.5, 6.5, 5.5} { \fill (-5.2,\y) circle (2pt); }

  % Alimentação
  \node[anchor=west, font=\small, align=left] at (0,2.5) {Conector\\[-2pt]de energia};
  \node[anchor=north, font=\small] at (1.5,0.6) {GND};
  \draw (-6.5,0) to[battery1, l=$9\,\mathrm{V}$] (-6.5,2.5)
        to[nos, l=chave] (-3.5,2.5) -- (0,2.5);
  \fill (-2,2.5) circle (2pt);
  \draw (-2,2.5) to[led, l_=\small vermelho] (-2,0.6) to[R, l=$1\,\mathrm{k\Omega}$] (-2,-1) -- (-2,-1.3);
  \draw (-6.5,0) -- (-6.5,-1.3) -- (1.5,-1.3) -- (1.5,0);
  \fill (-2,-1.3) circle (2pt);
  \draw (1.5,-1.3) node[ground]{};

  % Display OLED: alimentado por 5V e GND, dados em A4 (SDA) e A5 (SCL)
  \draw[thick] (0.1,10.6) rectangle (2.9,11.6);
  \node[font=\small\bfseries] at (1.5,11.1) {Display OLED};
  \foreach \x/\p in {0.4/GND, 1.1/5V, 1.9/A4, 2.6/A5} {
    \draw (\x,9) -- (\x,10.6);
    \node[anchor=south, font=\scriptsize] at (\x,9) {\p};
  }
\end{circuitikz}
}
	\caption{Esquema elétrico do LUXIF.}
	\label{fig:esquema}
\end{figure}

\subsection*{Versão com Wi-Fi}

Na versão com ESP32, o circuito de medição é o mesmo, com os pinos GPIO 25, 26, 27 e 14 para os resistores de referência, GPIO 16 a 19 para os LEDs e GPIO 34 para a leitura. O conversor analógico-digital do ESP32 tem 12 bits, e a leitura $N$ vai de 0 a 4095, com os limites de troca de faixa em 400 e 3700. As leituras são enviadas como texto JSON, no formato \texttt{\{"lux":523.4,"percentage":100.0,"faixa":1\}}, por dois caminhos: pela internet, com o protocolo MQTT (OASIS, 2014), em que o ESP32 publica as medições em um canal de um servidor e a página web assina esse canal; e pela rede local, com um servidor WebSocket (FETTE; MELNIKOV, 2011) acessado em \texttt{ws://IP/ws}, em que o navegador se conecta diretamente ao ESP32. O servidor WebSocket foi desenvolvido em ESP-IDF e usa o mesmo formato de mensagem, de modo que a mesma página exibe os dados vindos de qualquer um dos firmwares.

\subsection*{Conversão para lux}

A resistência de um LDR diminui com a iluminância segundo uma lei de potência (SENBA, [s.d.]):

\begin{equation}
	R_{LDR} = A \times L^{-\gamma}
	\label{eq:ldr}
\end{equation}

\noindent em que $L$ é a iluminância, em lux; $A$ é a resistência do LDR em 1~lux; e $\gamma$ indica a inclinação da curva. Foram adotados os valores típicos $A = 75$~k$\Omega$ e $\gamma = 0{,}7$ (Figura~\ref{fig:ldr}). Invertendo a Equação~\ref{eq:ldr}, o microcontrolador calcula a iluminância:

\begin{equation}
	L = \left( \frac{A}{R_{LDR}} \right)^{1/\gamma}
	\label{eq:lux}
\end{equation}

\begin{figure}[htbp]
	\centering
	\begin{tikzpicture}
  \begin{loglogaxis}[
    width=12cm, height=7cm,
    xmin=1, xmax=10000, ymin=100, ymax=100000,
    xlabel={Iluminância (lux)},
    ylabel={Resistência do LDR ($\Omega$)},
    grid=both, minor grid style={gray!15}, major grid style={gray!35},
  ]
    \addplot[domain=1:10000, samples=60, very thick, blue!70!black] {75000*x^(-0.7)};
  \end{loglogaxis}
\end{tikzpicture}

	\caption{Curva teórica da resistência do LDR em função da iluminância, em escala logarítmica.}
	\label{fig:ldr}
\end{figure}

\section*{RESULTADOS E DISCUSSÃO}

A Figura~\ref{fig:divisor} apresenta a tensão em A0 em função da resistência do LDR para cada resistor de referência, calculada pela Equação~\ref{eq:divisor}. A região sombreada corresponde à janela de leitura útil, entre 0,5 e 4,5~V. Cada curva é mais inclinada, e portanto mais sensível, quando a resistência do LDR é próxima de $R_n$; longe desse ponto, a curva fica plana e variações de luz quase não alteram a tensão. A troca automática mantém a medição sempre na parte inclinada de alguma das curvas.

\begin{figure}[htbp]
	\centering
	\begin{tikzpicture}
  \begin{semilogxaxis}[
    width=12cm, height=7cm,
    xmin=100, xmax=10000000, ymin=0, ymax=5,
    xlabel={Resistência do LDR ($\Omega$)},
    ylabel={Tensão em A0 (V)},
    grid=both, minor grid style={gray!15}, major grid style={gray!35},
    legend pos=south east, legend cell align=left,
  ]
    \addplot[fill=green!25, fill opacity=0.5, draw=none, forget plot] coordinates {(100,0.49) (10000000,0.49) (10000000,4.5) (100,4.5)} -- cycle;
    \addplot[domain=100:10000000, samples=120, very thick, blue!70!black] {5*x/(1000+x)};
    \addplot[domain=100:10000000, samples=120, very thick, orange!80!black] {5*x/(10000+x)};
    \addplot[domain=100:10000000, samples=120, very thick, green!50!black] {5*x/(100000+x)};
    \addplot[domain=100:10000000, samples=120, very thick, red!70!black] {5*x/(470000+x)};
    \legend{$R_n = 1\,\mathrm{k\Omega}$, $R_n = 10\,\mathrm{k\Omega}$, $R_n = 100\,\mathrm{k\Omega}$, $R_n = 470\,\mathrm{k\Omega}$}
  \end{semilogxaxis}
\end{tikzpicture}

	\caption{Tensão em A0 em função da resistência do LDR para os quatro resistores de referência. A região sombreada é a janela de leitura útil.}
	\label{fig:divisor}
\end{figure}

A Tabela~\ref{tab:faixas} resume a cobertura de cada faixa, obtida aplicando os limites de leitura (100 e 920) à Equação~\ref{eq:resistencia}. Cada faixa cobre resistências de cerca de $R_n/9$ a $9 \times R_n$. Como os resistores vizinhos diferem por um fator de 10 ou menos, as faixas se sobrepõem, o que cria uma histerese natural: quando a luz está perto de um limite, o instrumento não fica alternando entre duas faixas.

\begin{table}[H]
	\centering
	\caption{Cobertura teórica de cada faixa.}
	\label{tab:faixas}
	\small
	\begin{tabular}{cccc}
		\toprule
		Faixa & $R_n$ & $R_{LDR}$ coberta & Iluminância estimada \\
		\midrule
		1 & 1 k$\Omega$ & 108 $\Omega$ a 8,9 k$\Omega$ & 21 a 11\,000 lux \\
		2 & 10 k$\Omega$ & 1,1 k$\Omega$ a 89 k$\Omega$ & 0,8 a 425 lux \\
		3 & 100 k$\Omega$ & 11 k$\Omega$ a 893 k$\Omega$ & 0,03 a 16 lux \\
		4 & 470 k$\Omega$ & 51 k$\Omega$ a 4,2 M$\Omega$ & 0,003 a 1,7 lux \\
		\bottomrule
	\end{tabular}
\end{table}

As iluminâncias da Tabela~\ref{tab:faixas} dependem dos valores de $A$ e $\gamma$ adotados, que são os valores típicos da folha de dados e podem ser ajustados no código para cada sensor. A faixa de trabalho considerada para o LUXIF é de 1 a 1000~lux, na qual o sensor tem resposta mais previsível.

A ligação dos LEDs azuis fora do divisor e a do LED vermelho em paralelo com a alimentação evitam interferência na medição e sobrecarga no circuito, e utilizam apenas saídas livres do Arduino e resistores comuns.

\section*{CONCLUSÕES}

O presente trabalho apresentou o hardware de um luxímetro de baixo custo baseado em um LDR, em quatro resistores de referência e em um Arduino Uno. A análise teórica do divisor de tensão mostrou que a seleção automática entre os resistores mantém a leitura sempre na região de maior sensibilidade, ampliando a faixa dinâmica do instrumento em comparação com um resistor fixo.

Os LEDs indicadores, ligados a saídas próprias, não interferem na medição, e a alimentação por bateria de 9~V, com chave e LED indicador, torna o instrumento portátil. A versão com ESP32 transmite as leituras por Wi-Fi, pela internet (MQTT) ou pela rede local (WebSocket), para uma página web que exibe o valor em tempo real, o que torna o LUXIF adequado a experimentos de óptica e ao acompanhamento da iluminação natural.

\section*{CONTRIBUIÇÕES DOS AUTORES}

E.L.M.\textsuperscript{1} contribuiu como principal gerente e organizador do projeto, planejando, coordenando e conduzindo as pesquisas, refinando a base teórica e formulando os objetivos de estudo. H.F.D.\textsuperscript{3} e R.S.U.R.\textsuperscript{4} atuaram como principais investigadores e analistas formais, providenciando os dados e as ferramentas matemáticas, visuais e linguísticas, além de atuarem, de forma secundária, como supervisores e revisores. E.O.A.\textsuperscript{2} atuou como principal editora e revisora de dados, além de validar as informações obtidas e as ferramentas matemáticas utilizadas. V.M.P.\textsuperscript{5} desenvolveu o servidor WebSocket em ESP-IDF para o ESP32, usado na comunicação sem fio das leituras com a página web.

\section*{AGRADECIMENTOS}

Agradecimentos aos professores Dennis Toufen, Maurício Capelas, Hamilton Piva e Wilson Carlos por sua coordenação, ensino, condução e direção, que nos permitiram desenvolver e evoluir tanto este projeto e pesquisa como nossas vidas acadêmicas.

\section*{REFERÊNCIAS}

\begin{flushleft}
\small
ARDUINO. \textbf{Arduino Uno Rev3}: documentação técnica. [S.l.]: Arduino, [s.d.]. Disponível em: \url{https://docs.arduino.cc/hardware/uno-rev3/}.

\vspace{6pt}
HALLIDAY, David; RESNICK, Robert; WALKER, Jearl. \textbf{Fundamentos de Física}: Óptica e Física Moderna. 10. ed. Rio de Janeiro: LTC, 2016.

\vspace{6pt}
FETTE, Ian; MELNIKOV, Alexey. \textbf{RFC 6455}: The WebSocket Protocol. Internet Engineering Task Force, 2011. Disponível em: \url{https://www.rfc-editor.org/rfc/rfc6455}.

\vspace{6pt}
OASIS. \textbf{MQTT Version 3.1.1}. OASIS Standard, 2014. Disponível em: \url{https://docs.oasis-open.org/mqtt/mqtt/v3.1.1/mqtt-v3.1.1.html}.

\vspace{6pt}
SCHERZ, Paul; MONK, Simon. \textbf{Practical Electronics for Inventors}. 3. ed. New York: McGraw-Hill Education TAB, 2013.

\vspace{6pt}
SENBA OPTOELECTRONICS CO., LTD. \textbf{GL55 Series Photoresistor Datasheet}. Shenzhen: Senba Optoelectronics, [s.d.]. Disponível em: \url{https://www.digikey.com/en/htmldatasheets/production/2014143/0/0/1/gl5528.html}.
\end{flushleft}

\end{document}
